% Source: https://artofproblemsolving.com/community/c3h2412150_is_this_called_fun
% Stable identity: aops-topic-2412150-q01
\textbf{The numbers $0$ to $1999$.} Write them all with four digits, $0000$ to $1999$; the added leading zeros do not affect the count of $1$s and $2$s. In each of the last three places, every digit appears equally often: $\dfrac{2000}{10}=200$ times. So each of the digits $1$ and $2$ appears $3\cdot200=600$ times there. In the thousands place, $1$ appears $1000$ times (in $1000$ to $1999$) and $2$ never. So far: $1600$ ones and $600$ twos.

\textbf{The numbers $2000$ to $2021$.} The thousands digit is $2$ for all $22$ numbers, and the hundreds digit is $0$. The last two digits run through $00$ to $21$: the digit $1$ appears in $01,11,21$ (units) and $10,\dots,19$ (tens), $13$ times; the digit $2$ appears in $02,12$ (units) and $20,21$ (tens), $4$ times. So: $13$ ones and $22+4=26$ twos.

In total there are $1613$ ones and $626$ twos, and the difference is $1613-626=\boxed{987}$.
